\section*{Capítulo \ref{ch:b2:wittig} --- Wittig e outras reações que formam C=C}

\begin{solution}{exo:b2:wittig:1}
\ce{Ph3P + CH3CH2Br -> [Ph3PCH2CH3]+ + Br-} (substituição bimolecular, aquecimento em um solvente), depois
butil-lítio em éter anidro ou THF sob gás inerte: \ce{[Ph3PCH2CH3]+ + BuLi -> Ph3P=CHCH3 + BuH + Li+}.
\end{solution}

\begin{solution}{exo:b2:wittig:2}
O ilídeo \ce{Ph3P=CH2} e o benzaldeído dão o fenileteno (estireno), \ce{C6H5CH=CH2}, e o
óxido de trifenilfosfina.
\end{solution}

\begin{solution}{exo:b2:wittig:3}
\ce{Ph3P=CHCOOEt} e \ce{Ph3P=CHCOCH3}: estabilizados. \ce{Ph3P=CHCH3}: não estabilizado, o que
exige butil-lítio (ou hidreto de sódio, amideto de sódio). \ce{Ph3P=CHC6H5}: estabilizado por arila,
preparado na prática com um alcóxido ou hidróxido, e que dá misturas de geometrias.
\end{solution}

\begin{solution}{exo:b2:wittig:4}
(\textit{E})-pent-2-enoato de etila, \ce{CH3CH2CH=CHCOOEt}, com fosfato de dietila e sódio.
\end{solution}

\begin{solution}{exo:b2:wittig:5}
Wittig: o ilídeo \ce{Ph3P=CH2}, preparado a partir do iodeto de metiltrifenilfosfônio e do butil-lítio,
dá com a ciclo-hexanona apenas o metilidenociclo-hexano, com a ligação dupla fora do anel. A
desidratação do 1-metilciclo-hexan-1-ol passa por um carbocátion terciário e perde o hidrogênio que
dá o alceno mais substituído (Zaitsev): sobretudo 1-metilciclo-hexeno, o isômero errado.
\end{solution}

\begin{solution}{exo:b2:wittig:6}
Os dois cortes são iguais: benzaldeído e o ilídeo benzilideno \ce{Ph3P=CHC6H5}. Esse ilídeo é apenas
estabilizado por arila e dá uma mistura (\textit{E})/(\textit{Z}). Para o (\textit{E}) puro, use a
reação HWE do benzilfosfonato de dietila com o benzaldeído.
\end{solution}

\begin{solution}{exo:b2:wittig:7}
Com \ce{Ph3P=CHCH3}: (\textit{Z})-pent-2-eno, \ce{CH3CH2CH=CHCH3}, sobretudo (\textit{Z}). Com
\ce{Ph3P=CHCOOEt}: (\textit{E})-pent-2-enoato de etila, sobretudo (\textit{E}).
\end{solution}

\begin{solution}{exo:b2:wittig:8}
\ce{Ph3P=CH2 + C6H10O -> C7H12 + Ph3PO}. Massas: metilidenociclo-hexano \qty{96.173}{g/mol},
óxido de trifenilfosfina \qty{278.291}{g/mol}; economia atômica $96.173/(96.173 + 278.291) =
96.173/374.464 \approx 25.7~\%$. Três quartos da massa dos reagentes saem como óxido de fosfina.
% ledger: aw:C, aw:H, aw:O, aw:P
\end{solution}

\begin{solution}{exo:b2:wittig:9}
O fósforo, nucleofílico, desloca o brometo:
\begin{equation*}
\ce{P(OEt)3 + BrCH2COOEt -> [(EtO)3PCH2COOEt]+ + Br-}.
\end{equation*}
O brometo
ataca então um carbono de um grupo etila do íon fosfônio (substituição bimolecular), liberando
o bromoetano e formando a ligação \ce{P=O}:
\begin{equation*}
\ce{[(EtO)3PCH2COOEt]+ + Br- -> (EtO)2P(O)CH2COOEt + EtBr}.
\end{equation*}
\end{solution}

\begin{solution}{exo:b2:wittig:10}
\ce{OHC-CH=CH-CHO + 2Ph3P=CHC6H5 -> C6H5CH=CH-CH=CH-CH=CHC6H5 + 2Ph3PO}: dois equivalentes do
ilídeo benzilideno (a partir do cloreto de benziltrifenilfosfônio e de uma base), um para cada grupo
aldeído. As novas ligações duplas se formam como misturas de geometrias; o isômero todo-(\textit{E}),
o mais estável, é isolado por cristalização, ou a mistura é isomerizada.
\end{solution}

\begin{solution}{exo:b2:wittig:11}
\omterm{def:b2:enolates-aldol:aldol}{Aldólica}: benzaldeído com propanona em excesso e hidróxido de sódio (\cref{ch:b2:enolates-aldol});
reagentes baratos, água como subproduto, produto (\textit{E}), mas a condensação pode ocorrer duas
vezes (dibenzalacetona) se a propanona não estiver em excesso. Wittig: benzaldeído com o \omterm{def:b2:wittig:ylide}{ilídeo
estabilizado} \ce{Ph3P=CHCOCH3}; (\textit{E}), um só produto, sem dupla condensação, mas o ilídeo é caro
e o óxido de trifenilfosfina precisa ser removido. A \omterm{def:b2:enolates-aldol:aldol}{aldólica} é aqui a melhor via.
\end{solution}

\begin{solution}{exo:b2:wittig:12}
Wittig: butanal com \ce{Ph3P=CHCH2CH2CH3} (a partir do 1-bromobutano, depois butil-lítio), sem sais:
sobretudo (\textit{Z}); subproduto, o óxido de trifenilfosfina, e um pouco de isômero (\textit{E}).
Alcino: oct-4-ino (a partir do etino por duas alquilações com 1-bromopropano e amideto de sódio)
hidrogenado sobre um catalisador de paládio envenenado (\cref{prop:b2:alkene-redox:syn-hydrogenation}):
só (\textit{Z}); os subprodutos são o brometo de sódio e a amônia das alquilações.
\end{solution}

\begin{solution}{pb:b2:wittig:1}
\textbf{1.} \ce{C23H46}, isto é, \ce{C_nH_{2n}} com $n = 23$: uma ligação dupla (ou um anel), aqui uma
\ce{C=C}.
\textbf{2.} A ligação dupla \ce{C9=C10}.
\textbf{3.} O nonanal \ce{CH3(CH2)7CHO} com o ilídeo de um 1-halotetradecano; ou o tetradecanal
\ce{CH3(CH2)12CHO} com o ilídeo de um 1-halononano.
\textbf{4.} O 1-bromotetradecano, \ce{CH3(CH2)13Br}.
\textbf{5.} A substituição bimolecular pela volumosa trifenilfosfina é rápida em um carbono primário
desimpedido e não compete com a eliminação.
\textbf{6.} \ce{Ph3P + C14H29Br -> [Ph3PC14H29]+ + Br-}, brometo de tetradeciltrifenilfosfônio.
\textbf{7.} Butil-lítio (ou hidreto de sódio, amideto de sódio); o hidróxido é uma base fraca demais
para um ilídeo não estabilizado, e a água que ele traz protonaria o ilídeo.
\textbf{8.} Não, ele leva apenas um grupo alquila: sobretudo (\textit{Z}) (\cref{prop:b2:wittig:stereo}).
\textbf{9.} Um anel de quatro membros \ce{P-C-C-O}: o fósforo ligado ao carbono que leva a cadeia
\ce{C13H27}, esse carbono ao que leva a cadeia \ce{C8H17}, e este ao oxigênio, ligado ao
fósforo.
\textbf{10.} O butil-lítio e o ilídeo reagem com a água e o oxigênio.
\textbf{11.} A ligação dupla une o antigo carbono carbonílico do nonanal (C9 do produto) ao antigo
carbono do ilídeo (C10), e nenhuma outra ligação se desloca (\cref{prop:b2:wittig:regiospecific}).
\textbf{12.} Um carbocátion em C9, que pode perder um hidrogênio de C8 ou de C10 e pode migrar: uma
mistura de tricos-8-eno e tricos-9-eno, cada um como (\textit{E}) e (\textit{Z}), sobretudo (\textit{E}).
\textbf{13.} \ce{Ph3P=CHC13H27 + C8H17CHO -> C8H17CH=CHC13H27 + Ph3PO}.
\textbf{14.} Alvo \ce{C23H46}: \qty{322.621}{g/mol}; nonanal \ce{C9H18O}: \qty{142.242}{g/mol}; sal
\ce{C32H44BrP}: \qty{539.582}{g/mol}; \ce{Ph3PO}, \ce{C18H15OP}: \qty{278.291}{g/mol}.
\textbf{15.} O ilídeo \ce{C32H43P} pesa $539.582 - 80.912 = \qty{458.670}{g/mol}$; com o nonanal,
\qty{600.912}{g/mol} de reagentes para \qty{322.621}{g/mol} de produto: $322.621/600.912 \approx
53.7~\%$.
\textbf{16.} Pela polaridade: o hidrocarboneto apolar atravessa uma coluna curta de sílica com um
eluente apolar enquanto o óxido fica retido; ou o óxido cristaliza em um solvente apolar frio.
\textbf{17.} A ligação \ce{P=O} muito forte formada supera a ligação \ce{P-C} rompida
(\cref{prop:b2:wittig:driving}).
\textbf{18.} Ela exige um \omterm{def:b2:wittig:hwe}{fosfonato} com um grupo retirador de elétrons no carbono que reage, ausente
de uma simples cadeia alquila, e dá o alceno (\textit{E}).
\textbf{19.} O tricos-9-ino, \ce{CH3(CH2)7C#C(CH2)12CH3}, com hidrogênio sobre um catalisador de paládio
envenenado (paládio sobre carbonato de cálcio com um sal de chumbo e quinolina).
\textbf{20.} O dec-1-ino, \ce{CH3(CH2)7C#CH}, desprotonado pelo amideto de sódio, depois o 1-bromotridecano
\ce{CH3(CH2)12Br}: $8 + 2 + 13 = 23$ carbonos, a ligação tripla entre C9 e C10.
\textbf{21.} Uma hidrogenação adicional daria o tricosano, sem ligação dupla; o catalisador envenenado
liga o alcino muito mais fortemente que o alceno, que deixa a superfície antes de uma segunda adição.
\textbf{22.} Ambas levam duas etapas a partir de peças disponíveis. A via do alcino dá o isômero
(\textit{Z}) de modo limpo e deixa apenas sais; a via de Wittig dá um pouco de isômero (\textit{E}) e
uma massa de óxido de trifenilfosfina.
\textbf{23.} 100~\%: \ce{C23H44 + H2 -> C23H46}, todos os átomos terminam no produto.
\textbf{24.} $278.291/322.621 = \boldsymbol{\approx \qty{0.86}{g}}$ de óxido de trifenilfosfina por
grama de feromônio.
\end{solution}
